% Generated draft: requires mathematical and source review.
\documentclass[11pt]{article}
\usepackage[margin=25mm]{geometry}
\usepackage{fontspec}
\setmainfont{Latin Modern Roman}
\usepackage{amsmath,amssymb,mathtools,unicode-math}
\setmathfont{Latin Modern Math}
\usepackage{enumitem,needspace,xurl,hyperref}
\setlength{\parindent}{0pt}
\setlength{\parskip}{5pt}
\setlength{\emergencystretch}{3em}
\newcommand{\N}{\mathbb{N}}
\newcommand{\Z}{\mathbb{Z}}
\newcommand{\Q}{\mathbb{Q}}
\newcommand{\R}{\mathbb{R}}
\newcommand{\C}{\mathbb{C}}
\newcommand{\F}{\mathbb{F}}
\newcommand{\A}{\mathbb{A}}
\newcommand{\PP}{\mathbb{P}}
\newcommand{\E}{\mathbb{E}}
\newcommand{\eps}{\varepsilon}
\newcommand{\dd}{\,\mathrm{d}}
\begin{document}
\section*{UnsolvedMath ID: 701050}
\textbf{Source problem code:} catalogue entry
\subsection{Definitions and mathematical statement}
Let \(G\) be a finite group. For \(g\in G\), define its conjugacy class and centralizer by
\[
g^G=\{x^{-1}gx:x\in G\},\qquad C_G(g)=\{x\in G:xg=gx\}.
\]
The size of this conjugacy class is \(|g^G|=[G:C_G(g)]\). For a positive integer \(m\), let \(\pi(m)\) be the set of distinct prime divisors of \(m\), with \(\pi(1)=\varnothing\). Define
\[
\rho(G)=\bigcup_{g\in G}\pi(|g^G|),\qquad
\sigma(G)=\max_{g\in G}|\pi(|g^G|)|.
\]
Thus primes are counted without multiplicity, and the source's phrase ``order of a conjugacy class'' is understood as its cardinality. In particular, \(\rho(G)=\varnothing\) and \(\sigma(G)=0\) when \(G\) is abelian, including when \(G\) is trivial.

\textbf{Problem (catalogue entry; permanent catalogue entry).} Is it true that
\[
\forall\,G\text{ finite},\qquad |\rho(G)|\leq 3\sigma(G)?
\]
The supplied source describes this as a well-known problem, proposed by A. Moret\'{o} in 1999.

\textbf{Supplied background qualifications.} The source attributes to C. Casolo and S. Dolfi (1996) a family of finite groups \((G_n)\) for which the ratios are defined and
\[
\lim_{n\to\infty}\frac{|\rho(G_n)|}{\sigma(G_n)}=3.
\]
Accordingly, the source states that the coefficient \(3\) cannot be reduced in a universal bound of the displayed form. Its comment of 2009 reports a quadratic bound due to A. Moret\'{o} (2005) and a linear bound due to C. Casolo and S. Dolfi (2007). No constants or precise formulations of those earlier bounds are supplied. These background assertions and the registry's open status are supplied source assessments, not independently verified mathematical or literature conclusions.
\subsection{Short English statement}
For every finite group, must the total number of distinct primes occurring in its conjugacy class sizes be at most three times the largest number of distinct primes occurring in any one class size? The supplied source reports examples approaching the factor \(3\), and a 2009 comment reports earlier quadratic and linear bounds.
\subsection{Sources}
\begin{itemize}
\item Local UnsolvedMath registry: supplied complete source record, permanent ID \(701050\), exact source code \texttt{catalogue entry}, titled \emph{Kourovka Notebook Problem 14.65}. The record labels the problem open and was updated on 2026-10-16; this label is supplied metadata, not an independently verified status assessment. Upstream revision: \texttt{d371ee917ac9647fbb06cd6fc529deec6f3dda82}.
\item \emph{The Kourovka Notebook (Unsolved Problems in Group Theory)}, original problem 14.65, 14th issue (1999); supplied background also identifies the 21st issue edition (2026), arXiv:1401.0300v48. The supplied collection description names E. I. Khukhro and V. D. Mazurov as editors. Proposed by A. Moret\'{o}; described in the statement as a well-known problem.
\item C. Casolo and S. Dolfi, \emph{Rend. Sem. Mat. Univ. Padova} \textbf{96} (1996), 121--130. Supplied attribution for the family whose ratios tend to \(3\).
\item A. Moret\'{o}, \emph{Int. Math. Res. Not.} \textbf{2005}, no. 54, 3375--3383. The source's comment of 2009 attributes a quadratic bound to this reference.
\item C. Casolo and S. Dolfi, \emph{J. Group Theory} \textbf{10}, no. 5 (2007), 571--583. The source's comment of 2009 attributes a linear bound to this reference.
\end{itemize}
\end{document}
